A semaphore is an important synchronisation tool.  The idea (and name) is
inspired by railway signalling, where a flag is raised to indicate that track
ahead can be entered, or lowered to indicate that the track ahead is occupied.
A train arriving at the section of track waits until the flag is up, enters
the section, and the flag is lowered again.  When the train leaves the section
of track, the flag is raised. 

A semaphore 
%% has a boolean variable, that indicates whether the flag is down.  It
has two operations:
%
\begin{description}
\item[\scalastyle{down()}:]
Wait until the flag is in the up position, and then put the flag down and
procced;

\item[\scalastyle{up()}:]
Put the flag up, to allow another thread to proceed.
\end{description}
%
The operations \SCALA{down} and \SCALA{up} are sometimes called \SCALA{P} and
\SCALA{V} (the initial letters of the Dutch words for the operations),
{\scalastyle wait} and \SCALA{signal}, or \SCALA{acquire} and
\SCALA{release}.

Figure~\ref{fig:semaphore-implementation} gives a straightforward
implementation of a semaphore using a JVM monitor.  The variable |isUp|
records whether the flag is currently in the up position.

%%%%%

\begin{figure}
\begin{scala}
/** A binary semaphore.
  * @param £isUp£ is the semaphore initially in the up state? */
class Semaphore(private var isUp: Boolean){
  def down() = synchronized{
    while(!isUp) wait()
    isUp = false
  }

  def up() = synchronized{
    require(!isUp); isUp = true; notify()
  }
}
\end{scala}
\caption{An implementation of a binary semaphore.}
\label{fig:semaphore-implementation}
\end{figure}

Note that the |up| operation requires that the semaphore is not already up.
This is a slightly non-standard design decisions: most implementations treat
the operation as a no-op in this case.  However, my experience is that trying
to put the semaphore up when it is already up is nearly always an error: a
signal gets lost.  I think it is better for the implementation to throw an
exception in this case, so that the error can be detected and removed.

\begin{instruction}
Make sure you understand how the implementation matches the informal
description given earlier.
\end{instruction}

In fact, many operating system kernels provide implementations of semaphores.
Those semaphores can then be used to implement other concurrency support, such
as monitors, within a programming language.

As with other concurrency mechanisms, semaphores can be used both for
providing mutual exclusion, which we will study in the next section, and for
waiting and signalling, which we will see subsequently.

%%%%%%%%%%%%%%%%%%%%%%%%%%%%%%%%%%%%%%%%%%%%%%%%%%%%%%%%%%%%

\section{Using a Semaphore for Mutual Exclusion}

One use of semaphores is to provide mutual exclusion.  

Such a semaphore is created with
\begin{scala}
  val mutex = new Semaphore(true)
\end{scala}
or
\begin{scala}
  val mutex = new MutexSemaphore
\end{scala}
%
using the definition
%
\begin{scala}
  class MutexSemaphore extends Semaphore(true)
\end{scala}
%
A thread entering the critical region executes
\begin{scala}
  mutex.down()
\end{scala}
%
This blocks if the semaphore is already down, i.e.~if another thread is in the
critical region.  In the initial state, the operation will not block. 
%
A thread leaving the critical region executes
\begin{scala}
  mutex.up()
\end{scala}
%
This unblocks a waiting thread (if there is one).

Figure~\ref{fig:semaphore-counter} gives an implementation of a counter that
uses a semaphore to provide mutual exclusion on the private variable~|x|.
Each operation initially performs a |down| on the semaphore, and performs an
|up| at the end.  Note that the |get| operation has to copy the value of~|x|
into a thread-local variable before raising the semaphore.

%%%%%

\begin{figure}
\begin{scala}
object Counter{
  private var x = 0
  
  private var mutex = new MutexSemaphore
    
  def inc() = { mutex.down(); x = x+1; mutex.up() }
  def dec() = { mutex.down(); x = x-1; mutex.up() }
  def get = { mutex.down(); val result = x; mutex.up(); result }
}
\end{scala}
\caption{An implementation of a counter using a semaphore for mutual
  exclusion.}
\label{fig:semaphore-counter}
\end{figure}

%%%%%

This technique can be used to provide mutual exclusion on arbitrary objects,
much like a |Lock|.

As with most concurrency primitives, it is possible to misuse semaphores.  For
exampe, consider the following code, where semaphore \SCALA{sem1} is used to
protect resource \SCALA{res1}, and semaphore \SCALA{sem2} is used to protect
resource \SCALA{res2}.
%
\begin{scala}
  val sem1 = MutexSemaphore(); val sem2 = MutexSemaphore()
  def t1 = thread{ sem1.down(); ...; sem2.down(); ... ; sem2.up(); ...; sem1.up() }
  def t2 = thread{ sem2.down(); ...; sem1.down(); ... ; sem1.up(); ...; sem2.up() }
  run(t1 || t2)
\end{scala}
%
This can deadlock.  Thread~|t1| can perform  |sem1.down()|, and thread~|t2| can
perform~|sem2.down()|.  Then, when each tries to perform |down| on the other
semaphore, they are blocked indefinitely. 

